\documentclass[a4paper,12pt]{extarticle}
\usepackage{amsmath}
\usepackage{amssymb}
\usepackage{dsfont}
\usepackage[lmargin=71pt, tmargin=1.2in]{geometry}  %For centering solution box
% \chead{\hline} % Un-comment to draw line below header
\thispagestyle{empty}   %For removing header/footer from page 1
\usepackage{geometry}
\usepackage[T2A,T1]{fontenc}
\usepackage[utf8]{inputenc}
\usepackage[english,russian]{babel}
\usepackage{amsmath}
\usepackage{amsthm}
\usepackage{amssymb}
\usepackage{fancyhdr}
\usepackage{setspace}
\usepackage{graphicx}
\usepackage{colortbl}
\usepackage{tikz}
\usepackage{pgf}
\usepackage{subcaption}
\usepackage{listings}
\usepackage{indentfirst}

\usepackage{listings}
\usepackage{xcolor}

\usepackage{fvextra}

\begin{document}

\begingroup  
    \centering
    \LARGE Оптимальное инвестирование и риск-менеджмент\\
    \LARGE  Контрольная №2\\[0.5em]
    \large Задания (1, 1, 3, 1, 3)
    \large \today\\[0.5em]
    \large Детер Климентий\par
    \large Прикладная математика и информатика 2310\par
\endgroup
% \rule{\textwidth}{0.4pt}

\begin{center}
    \textbf{Задача 1(Вариант 1)}
\end{center}
При каких $a,b$ процесс $B^3_t +t + atB_t + bB_t$ является мартингалом относительно фильтрации $\mathcal{F}_t = \sigma(B_s, s\leq t)$, где $B_t$ — броуновское движение?
Чему равна $P(B_6 > 0,B_3 <0)$?
\begin{center}
    \textbf{Решение}
\end{center}
\begin{enumerate}
    \item Найдем дифференциал $dX_t$ по лемме Ито:
    $$df = \frac{\partial f}{\partial t}dt + \frac{\partial f}{\partial x}dB_t + \frac{1}{2}\frac{\partial^2 f}{\partial x^2}(dB_t)^2$$
    Производные:
    $$\frac{\partial f}{\partial t} = 1 + aB_t$$
    $$\frac{\partial f}{\partial x} = 3B_t^2 + at + b$$
    $$\frac{\partial^2 f}{\partial x^2} = 6B_t$$
    Подставляем в формулу (учитывая $(dB_t)^2 = dt$):
    $$dX_t = (1 + aB_t)dt + (3B_t^2 + at + b)dB_t + \frac{1}{2}(6B_t)dt$$
    $$dX_t = (1 + (a+3)B_t)dt + (3B_t^2 + at + b)dB_t$$

    \item Для мартингальности дрейф должен быть тождественно равен нулю:
    $$1 + (a+3)B_t \equiv 0$$
    Это должно выполняться для любых $t$ и $B_t$. Поскольку свободный член $1 \neq 0$, уравнение не имеет решений, не зависящих от реализации случайного процесса $B_t$.

    \item Вычисление вероятности:
    Пусть $X = B_3$ и $Y = B_6 - B_3$. По свойствам винеровского процесса, $X$ и $Y$ — независимые случайные величины с распределением $\mathcal{N}(0, 3)$.
    Ищем $P(X + Y > 0, X < 0)$, что эквивалентно $P(Y > -X, X < 0)$.
    Вектор $(X, Y)$ — стандартный двумерный нормальный с независимыми компонентами, поэтому ротационно симметричен.
    В полярных координатах $X = r\cos\theta$, $Y = r\sin\theta$, угол $\theta \sim \mathrm{Uniform}[0, 2\pi)$.
    Условие $X < 0$ даёт $\theta \in (\pi/2,\, 3\pi/2)$; условие $Y > -X$, то есть $X + Y > 0$, даёт
    $\sin\theta + \cos\theta > 0 \iff \theta \in [0,\, 3\pi/4) \cup (7\pi/4,\, 2\pi)$.
    Пересечение: $\theta \in (\pi/2,\, 3\pi/4)$, длина интервала $3\pi/4 - \pi/2 = \pi/4$.
    Вероятность: $\dfrac{\pi/4}{2\pi} = \dfrac{1}{8}$.
\end{enumerate}

\textbf{Ответ:}
\begin{itemize}
    \item Таких $a,b$ не существует, процесс не может быть мартингалом из-за неустранимого дрейфа $dt$.
    \item $P(B_6 > 0, B_3 < 0) = \frac{1}{8} = 0.125$.
\end{itemize}

\begin{center}
    \textbf{Задача 2(Вариант 1)}
\end{center}
Пусть $\xi_1$,$\xi_2$ — н.о.р.с.в, имеющие равномерное распределение на
[a,b]. Пусть $\eta = \xi_1 - \xi_2$. Найти эксцесс и асимметрию у $\eta$, а также $VaR_{10\%}(\eta), ESF_{10\%}(\eta)$.
\begin{center}
    \textbf{Решение}
\end{center}

Пусть $c = b-a$. Величина $\eta$ имеет симметричное треугольное распределение на интервале $[-c, c]$.
Плотность: $f(x) = \frac{c - |x|}{c^2}$. Риск-метрики считаются для распределения потерь $L = -\eta$.
\begin{enumerate}
    \item \textbf{Асимметрия и Эксцесс:}
    $\mathbb{E}[\eta] = 0$. В силу симметрии распределения третий центральный момент равен нулю, следовательно, асимметрия равна 0.
    Дисперсия:
    $$Var(\eta) = Var(\xi_1) + Var(\xi_2) = \frac{c^2}{12} + \frac{c^2}{12} = \frac{c^2}{6}$$
    Четвертый момент:
    $$\mathbb{E}[\eta^4] = 2 \int_0^c x^4 \frac{c-x}{c^2} dx = \frac{2}{c^2} \left[ c\frac{c^5}{5} - \frac{c^6}{6} \right] = \frac{c^4}{15}$$
    Эксцесс:
    $$Kurtosis = \frac{\mathbb{E}[\eta^4]}{Var(\eta)^2} - 3 = \frac{c^4 / 15}{c^4 / 36} - 3 = 2.4 - 3 = -0.6$$

    \item \textbf{Value at Risk ($VaR_{10\%}$):}
    Условие $P(L > VaR) = 0.1 \implies P(\eta < -VaR) = 0.1$.
    Функция распределения (левый хвост) для $x \in [-c, 0]$: $F(x) = \frac{(x+c)^2}{2c^2}$.
    $$\frac{(-VaR+c)^2}{2c^2} = 0.1 \implies (-VaR+c)^2 = 0.2c^2$$
    $$VaR_{10\%} = c(1 - \sqrt{0.2}) = (b-a)\left(1 - \frac{1}{\sqrt{5}}\right)$$

    \item \textbf{Expected Shortfall ($ESF_{10\%}$):}
    $$ESF_{\alpha} = \frac{1}{\alpha} \int_0^{\alpha} VaR_p dp$$
    $$VaR_p = c(1 - \sqrt{2p})$$
    $$ESF_{0.1} = 10 \int_0^{0.1} c(1 - \sqrt{2p}) dp = 10c \left[ p - \frac{2(2p)^{3/2}}{3\cdot 2} \right]_0^{0.1}$$
    $$ESF_{0.1} = 10c \left( 0.1 - \frac{0.2\sqrt{0.2}}{3} \right) = (b-a)\left(1 - \frac{2}{3\sqrt{5}}\right)$$
\end{enumerate}

\textbf{Ответ:}
\begin{itemize}
    \item Асимметрия: $0$
    \item Эксцесс: $-0.6$
    \item $VaR_{10\%} = (b-a)\left(1 - \frac{1}{\sqrt{5}}\right) \approx 0.553(b-a)$
    \item $ESF_{10\%} = (b-a)\left(1 - \frac{2}{3\sqrt{5}}\right) \approx 0.702(b-a)$
\end{itemize}

\begin{center}
    \textbf{Задача 3(Вариант 3)}
\end{center}
Пусть $X_1$,$X_2$ — н.о.р.с.в., имеющие экспоненциальное распределение с параметром $\lambda$. Пусть $Y= X_1 + X_2$. Найти риск-вклад $X_1$ в Y, используя метрики: Componental Expected Shortfall с параметром 0,05, Incremental VaR.
\begin{center}
    \textbf{Решение}
\end{center}

Сумма $Y \sim Gamma(2, \lambda)$ с плотностью $f_Y(y) = \lambda^2 y e^{-\lambda y}$ и $F_Y(y) = 1 - e^{-\lambda y}(1+\lambda y)$.
Квантиль $y_{\alpha} = VaR_{0.05}(Y)$ находится из уравнения $F_Y(y) = 0.95$.
По определению: $CES_1 = \mathbb{E}[X_1 \mid Y \geq y_{\alpha}]$ и $IVaR_1 = \mathbb{E}[X_1 \mid Y = y_{\alpha}]$.

\textbf{Обоснование тождеств.}
По определению Expected Shortfall на уровне $\alpha$:
$$ES_{\alpha}(Y) = \mathbb{E}[Y \mid Y \geq y_{\alpha}]$$
Так как $Y = X_1 + X_2$, по линейности условного математического ожидания:
$$ES_{\alpha}(Y) = \mathbb{E}[X_1 + X_2 \mid Y \geq y_{\alpha}] = \mathbb{E}[X_1 \mid Y \geq y_{\alpha}] + \mathbb{E}[X_2 \mid Y \geq y_{\alpha}]$$
Поскольку $X_1$ и $X_2$ — независимые и одинаково распределённые, их роли в сумме $Y$ симметричны, поэтому
$$\mathbb{E}[X_1 \mid Y \geq y_{\alpha}] = \mathbb{E}[X_2 \mid Y \geq y_{\alpha}]$$
Отсюда $ES_{\alpha}(Y) = 2\mathbb{E}[X_1 \mid Y \geq y_{\alpha}]$, то есть
$$CES_1 = \frac{1}{2}ES_{\alpha}(Y)$$

Аналогично для Incremental VaR. Так как $VaR_{\alpha}(Y) = y_{\alpha}$, имеем
$$VaR_{\alpha}(Y) = \mathbb{E}[Y \mid Y = y_{\alpha}] = \mathbb{E}[X_1 \mid Y = y_{\alpha}] + \mathbb{E}[X_2 \mid Y = y_{\alpha}]$$
В силу той же симметрии $\mathbb{E}[X_1 \mid Y = y_{\alpha}] = \mathbb{E}[X_2 \mid Y = y_{\alpha}]$, откуда
$$IVaR_1 = \frac{1}{2}VaR_{\alpha}(Y)$$

\begin{enumerate}
    \item Решение уравнения квантили:
    Пусть $z = \lambda y_{\alpha}$. $e^{-z}(1+z) = 0.05$. Численное решение даёт $z \approx 4.74386$.
    $$VaR_{0.05}(Y) = \frac{4.74386}{\lambda}$$

    \item Численный расчёт:
    $$IVaR_1 = \frac{1}{2}VaR_{0.05}(Y) = \frac{2.37193}{\lambda}$$
    Общий $ESF_{0.05}(Y)$:
    $$ESF_{0.05}(Y) = \frac{1}{0.05} \int_{y_{\alpha}}^{\infty} y \lambda^2 y e^{-\lambda y} dy = \frac{20}{\lambda} [e^{-z}(z^2+2z+2)]$$
    Используя $e^{-z} = \frac{0.05}{1+z}$:
    $$ESF_{0.05}(Y) = \frac{1}{\lambda} \frac{z^2+2z+2}{z+1} \approx \frac{5.9179}{\lambda}$$
    $$CES_1 = \frac{1}{2} ESF_{0.05}(Y) \approx \frac{2.95895}{\lambda}$$
\end{enumerate}

\textbf{Ответ:}
\begin{itemize}
    \item $IVaR_1 \approx \frac{2.372}{\lambda}$
    \item $Componental Expected Shortfall (CES_1) \approx \frac{2.959}{\lambda}$
\end{itemize}

\begin{center}
    \textbf{Задача 4(Вариант 1)}
\end{center}
Из набора акций (AAPL, TSLA, IBM, AMZN) с 2015 года по конец 2024 года берете c yfinance нужную по счету акцию, соответствующую номеру Вашей фамилии по этой задачи и следующую за ней акцию. В предположении, что безрисковая $\%$ ставка $\dfrac{3.5\%}{12}$ в месяц постройте оптимальный портфель с помощью подхода Келли, а также на тестовой части выборки (2025 год) провести бэктестинг полученной стратегии, найдя норму доходность и Sharpe Ratio. Ограничение: стандартный Pyton + библиотека numpy + оптимизатор на Ваш выбор (если нужен) (cvxopt, например).
\begin{center}
    \textbf{Решение}
\end{center}


\begin{center}
    \textbf{Задача 5(Вариант 3)}
\end{center}
Пусть имеется безрисковый актив с $r^0 = 6$ и 2 рисковых с $\mathbb{E}(r^1) = 8$,$\sigma^1 = 5$,$\mathbb{E}(r^2) = 20$,$\sigma^2 = 20$,$corr(r^1,r^2) = 0$. Чему равна максимальная средняя прибыль и в каких пропорциях нужно вкладывать в активы, если $\sigma(r) \leq 15$ (короткие продажи разрешены).
\begin{center}
    \textbf{Решение}
\end{center}

Ковариационная матрица $\Sigma = diag(25, 400)$.\\
Целевая функция: $\mu_p = w_0 r^0 + w_1 \mu_1 + w_2 \mu_2$.\\
Ограничения: $w_0 + w_1 + w_2 = 1$ и $w_1^2 \sigma_1^2 + w_2^2 \sigma_2^2 \leq 225$.\\
\textbf{Множители Лагранжа}
\\Ограничение активно. Функция Лагранжа:
$$L = 6 + w_1(2) + w_2(14) - \lambda(25w_1^2 + 400w_2^2 - 225)$$
Условия ККТ:
$$\frac{\partial L}{\partial w_1} = 2 - 50\lambda w_1 = 0 \implies w_1 = \frac{1}{25\lambda}$$
$$\frac{\partial L}{\partial w_2} = 14 - 800\lambda w_2 = 0 \implies w_2 = \frac{7}{400\lambda}$$
Подставляем в ограничение риска:
$$25\left(\frac{1}{25\lambda}\right)^2 + 400\left(\frac{7}{400\lambda}\right)^2 = 225 \implies \lambda = \frac{\sqrt{65}}{300}$$
Подстановка $\lambda$ обратно дает те же самые значения $w_1$ и $w_2$.

\textbf{Ответ:}
\begin{itemize}
    \item Максимальная средняя прибыль: $6 + \frac{3\sqrt{260}}{4} \approx 18.093\%$
    \item Веса в рисковые активы: $w_1 = \frac{24}{\sqrt{260}} \approx 148.8\%$, $w_2 = \frac{10.5}{\sqrt{260}} \approx 65.1\%$
    \item Доля в безрисковом активе: $w_0 = 1 - \frac{34.5}{\sqrt{260}} \approx -114.0\%$ (кредитная позиция)
\end{itemize}
\end{document}